\section{Lesson 4: Crossing State Lines Fragments Districts and Hurts Compactness}
\label{sec:lesson4}

National redistricting without state lines (Paper E.3) is the purest
implementation of geometric optimization: treat the continental United States
as a single graph of approximately 73,000 census tracts and partition it into
435 compact contiguous districts using the METIS minimum-cut algorithm, without
any constraint that district boundaries respect state lines
\citep{dellaLibera2026e3}.  The design lesson from this experiment is
surprising and counterintuitive: E.3 finds that removing state boundaries
creates a \textbf{42.5\% compactness penalty} (PP 0.265 vs.\ 0.461 under the
50-state state-constrained baseline).
Geometric fragmentation of congressional districts across state lines
outweighs any optimization benefit.
This finding implies state boundaries are \emph{compactness-enabling} constraints,
not compactness-limiting ones.

\subsection{The Compactness Penalty}

The 50-state baseline achieves a mean Polsby--Popper score of 0.461.
The national redistricting experiment achieves only 0.265 --- a drop of
$\Delta\mathrm{PP} = -0.196$, or 42.5\% below the state-constrained baseline
\citep{dellaLibera2026e3}.

The penalty arises because districts that cross state lines must traverse
long, irregular boundaries where two states' census-tract grids meet at
misaligned angles.
State lines in the continental United States often run along rivers, mountain
ridges, and surveying meridians that themselves form compact geographic contours
within each state.
When a district crosses a state line, it must bridge two different geographic
coordinate systems and two different population density gradients, producing
elongated, irregular shapes that score poorly on the Polsby--Popper metric.
In this respect, state boundaries happen to coincide with the natural geometric
seams of the American landscape, and crossing them breaks compactness rather
than enabling it.

Approximately 60\% of districts are unchanged or minimally changed by removing
the state constraint: they would fall entirely within one state under any
reasonable compact-district drawing.  The compactness gain is concentrated in
the remaining 40\%, which gain an average of $\Delta\mathrm{PP} = +0.110$
--- a substantial improvement in the districts where the state constraint
was binding.

\subsection{The County Preservation Cost}

The compactness improvement comes at a significant cost in county preservation.
The baseline DIA preserves 38\% of counties whole (no tract splits across
district lines).  The national redistricting experiment preserves only 23\%
of counties whole --- a 40\% reduction in county-preservation rate.

This is counterintuitive at first: if the algorithm is optimizing for
compactness without state-line constraints, why does county preservation
decline?  The answer is that state boundaries often happen to coincide with
county boundaries --- state lines commonly run along county borders, especially
in the South and Midwest where counties were established contemporaneously with
state boundaries.  When the algorithm is free to cross state lines, it draws
districts that cross county lines in new places, fragmenting counties that
had been preserved under the state-constrained baseline.

The tradeoff between compactness and county preservation is thus partially
mediated by the state-boundary constraint: state boundaries impose compactness
costs in some places while incidentally protecting county integrity in others.
Removing the state constraint improves geometric compactness while degrading
county coherence.

\subsection{Political Fragmentation}

Beyond county preservation, national redistricting produces a specific type of
political fragmentation that does not arise in the state-constrained baseline:
districts that cross state lines represent constituents governed by different
state laws.

A district spanning the Kansas--Missouri border would represent residents of
two states with different income tax rates, different Medicaid eligibility
thresholds, different state legislative districts, different US Senate
representation, and different electoral calendars.  The House member from such
a district would face constituent questions --- about state programs, state
regulations, state courts --- that do not map cleanly onto her congressional
role.  She would also face a reelection challenge defined by two different
state primary calendars and two different state party structures.

Paper E.3 identifies 87 districts with $\geq 10\%$ cross-state population in the national map
(from E.3's sensitivity analysis), districts
that draw at least 10\% of their population from a different state than the
majority of their population \citep{dellaLibera2026e3}.  These 87 districts
represent a substantial fraction of the 435 total (20\%) and are concentrated
in the densely settled Northeast and the suburban Midwest.

\subsection{The Federalism Cost: A Precise Measurement}

The central contribution of E.3 to the design lessons is that it makes the
federalism cost of redistricting legible.  Before this experiment, the cost of
respecting state boundaries in redistricting was unmeasured; the conventional
wisdom was simply that state-based redistricting is constitutionally required
and therefore not worth analyzing as a design choice.

E.3 finds that removing state boundaries creates a 42.5\% compactness
\emph{penalty} (PP 0.265 vs.\ 0.461 under the state-constrained baseline).
This is the opposite of the intuition that unconstrained national optimization
would improve compactness.
State boundaries are compactness-enabling constraints: they confine districts
to coherent geographic regions whose census-tract adjacency graphs form compact
subgraphs.
Cross-state districts must span two different regional graph structures,
producing irregular shapes that score poorly on Polsby--Popper.

The implication for reform advocates is that the compactness argument for
national redistricting is not merely weak --- it is negative.  National
redistricting would produce geometrically worse districts than the
state-constrained approach.  The proportionality argument --- that national
redistricting might reduce the sorting-driven seat gap by drawing districts
that cross state-level sorting concentrations --- is potentially stronger but
is not supported by E.3's empirical results (see Lesson~1 above).
